\documentclass[11pt,a4paper]{article}
\usepackage[utf8]{inputenc}
\usepackage[T1]{fontenc}
\usepackage[spanish]{babel}
\usepackage{hyperref}
\usepackage{geometry}
\geometry{margin=2.5cm}
\usepackage{parskip}

\title{076 los alienigenas de las peliculas son humanos con otra piel}
\author{Alejandro Valenzuela Rivera\\\small Universidad Autónoma de Nuevo León}
\date{\today}
\begin{document}
\maketitle

\begin{quote}
\textbf{Sesgo declarado.} Este trabajo se sostiene desde el sesgo declarado de su autor, y no finge neutralidad. Lo que afirma se puede tumbar, y ese es el punto: el falsificador está escrito, no escondido.
\end{quote}

\section*{Resumen}
Un alien que siente como nosotros, se enoja como nosotros y quiere lo mismo que nosotros no es un alien: es un humano con otra piel. Las peliculas cambian la piel, el idioma y el nombre, y dejan intacto lo que hace que nos entendamos.

\section{Por qué se sostiene}
No es un defecto artistico: es un limite del que lo escribe, porque se escribe sobre lo que se puede leer, y lo que se puede leer es humano. Y por eso buscar aliens con forma de humano garantiza no encontrarlos. Lo que se hace es leer su mecanismo desde sus zapatos: su arena, su cuenta, su esquema y su borde. Si se puede correr sin el, tiene algo como la IA: un proceso, no una piel.

\section{De dónde viene}
El primer commit y la conversacion de la era.

\section{El caso}
Las peliculas de alien: humanos con otra piel.

\section{Método}
Una forma nueva que no es la vieja pintada de otro color.

\section{Límites, y qué lo tumba}
Una diferencia que solo es la forma de siempre con otra piel.

\section{Lo que queda abierto}
\textquestiondown{}Que hace un ser que no entiende la pregunta cuando le pido algo que no esta en su forma?

\section*{Cómo se leyó esta tesis}
Esto se escribió en el repositorio man-made-empathy, del que este candado es una pieza: candados/06 del caso/076 los alienigenas de las peliculas son humanos con otra piel.md. La máquina que lo compuso lo hizo desde este mismo archivo, y no agregó nada que no estuviera escrito. Si algo de esto es falso, se dice en el mismo sitio, y se corrige con fecha.

\end{document}
